\section{Entrenamiento e Implementación}

\subsection{Hiperparámetros}
\label{sec:hiperparam}

La Tabla~\ref{tab:hiperparam} resume los hiperparámetros
empleados para entrenar el Modelo~1 y su comparativa OkanNet.
Para garantizar una comparación científicamente limpia, se
mantienen idénticos en ambas arquitecturas; la única diferencia
controlada entre ellas es la presencia de
\textit{Batch Normalization}.

\begin{table}[!t]
\centering
\caption{Hiperparámetros del Modelo~1 (Baseline) y OkanNet.}
\label{tab:hiperparam}
\renewcommand{\arraystretch}{1.15}
\footnotesize
\begin{tabular}{@{}lll@{}}
\toprule
\textbf{Hiperparámetro} & \textbf{Baseline} & \textbf{OkanNet} \\
\midrule
Optimizador               & Adam                      & Adam                      \\
\quad $(\beta_1, \beta_2)$ & $(0{,}9, 0{,}999)$        & $(0{,}9, 0{,}999)$        \\
Learning rate inicial     & $1\times 10^{-3}$         & $1\times 10^{-3}$         \\
Scheduler                 & ReduceLROnPlateau         & ReduceLROnPlateau         \\
\quad factor / patience   & $0{,}5$ / $3$             & $0{,}5$ / $3$             \\
\quad $\eta_{\min}$       & $1\times 10^{-6}$         & $1\times 10^{-6}$         \\
Función de pérdida        & \texttt{CrossEntropyLoss} & \texttt{CrossEntropyLoss} \\
\quad con                 & \texttt{class\_weight}    & \texttt{class\_weight}    \\
\textit{Weight decay}     & $0$                       & $0$                       \\
\textit{Batch size}       & $32$                      & $32$                      \\
\textit{Epochs} máx.      & $40$                      & $40$                      \\
\textit{Early stopping}   & patience $=7$             & patience $=7$             \\
\quad criterio            & F1-macro val.             & F1-macro val.             \\
\textit{Split} train/val  & 80\,\%\,/\,20\,\%          & 80\,\%\,/\,20\,\%          \\
\textit{Seed}             & $42$                      & $42$                      \\
\bottomrule
\end{tabular}
\end{table}

El uso de \texttt{class\_weight} en la función de pérdida es un
diferencial del presente trabajo frente a las referencias
revisadas~\cite{okannet2026,gupta2023}, ninguna de las cuales
incorpora corrección explícita del desbalance de clases. Los
pesos se calculan con la fórmula \textit{balanced} de
\texttt{scikit-learn}, inversamente proporcionales a la
frecuencia de cada clase en \textit{Training}.

El tiempo de entrenamiento sobre Apple Silicon con backend MPS
fue de $7$\,min $44$\,s para el Baseline y $3$\,min $39$\,s
para OkanNet, este último interrumpido por \textit{early
stopping} en la época~$17$.

\subsubsection*{Hiperparámetros del Modelo~2}
\label{sec:hiperparam-m2}

Las dos variantes del Modelo~2 (DeepCNN-BN y DeepCNN-BN-GAP)
comparten un régimen de entrenamiento propio, distinto al del
Modelo~1, motivado por el fallo de OkanNet:

\begin{table*}[!t]
\centering
\caption{Hiperparámetros del Modelo~2 (DeepCNN-BN y
DeepCNN-BN-GAP). Idénticos en ambas variantes para aislar el
efecto del cabezal.}
\label{tab:hiperparam-m2}
\renewcommand{\arraystretch}{1.15}
\footnotesize
\begin{tabular}{@{}lll@{}}
\toprule
\textbf{Hiperparámetro} & \textbf{Modelo 2} & \textbf{vs Modelo 1} \\
\midrule
Optimizador               & AdamW                     & Adam                    \\
\quad $(\beta_1, \beta_2)$ & $(0{,}9, 0{,}999)$        & idéntico                \\
\textit{Learning rate}    & $5\times 10^{-4}$         & $10^{-3}$ (M1)          \\
\textit{Weight decay}     & $10^{-4}$                 & $0$                     \\
\textit{Batch size}       & $64$                      & $32$ (M1)               \\
Scheduler                 & ReduceLROnPlateau         & idéntico                \\
\quad factor / patience   & $0{,}5$ / $3$             & idéntico                \\
\quad $\eta_{\min}$       & $10^{-6}$                 & idéntico                \\
Función de pérdida        & \texttt{CrossEntropyLoss} & idéntico                \\
\quad con                 & \texttt{class\_weight}    & idéntico                \\
\textit{Epochs} máx.      & $50$                      & $40$ (M1)               \\
\textit{Early stopping}   & patience $=8$             & patience $=7$ (M1)      \\
\quad criterio            & F1-macro val.             & idéntico                \\
\textit{Split} train/val  & 80\,\%\,/\,20\,\%         & idéntico                \\
\textit{Seed}             & $42$                      & idéntico                \\
\bottomrule
\end{tabular}
\end{table*}

Tres ajustes distinguen al Modelo~2 del Modelo~1:

\begin{enumerate}
  \item \textbf{Optimizador AdamW en lugar de Adam.} La
        diferencia clave es que AdamW desacopla
        \textit{weight decay} del gradiente, evitando el
        sesgo que Adam clásico introduce al combinar
        \textit{L2 regularization} con momento
        adaptativo~\cite{loshchilov2017decoupled}. Esto
        habilita el uso de \textit{weight decay} no nulo
        ($10^{-4}$) como regularización adicional.
  \item \textbf{\textit{Learning rate} reducido y
        \textit{batch size} duplicado.} La combinación
        $\text{lr}=10^{-3}$ con $\text{batch}=32$
        empleada en el Modelo~1 generó inestabilidad en
        OkanNet (véase Tabla~\ref{tab:metricas-global}). En el
        Modelo~2 se baja \textit{lr} a $5\times 10^{-4}$ y se
        sube \textit{batch size} a $64$ para reducir la
        varianza de las estadísticas estimadas por
        \textit{mini-batch} en BN: con $\text{batch}=64$ la
        estimación de medias y varianzas es más estable, lo
        que evita el ``ruido amplificado'' que en OkanNet
        provocó la divergencia.
  \item \textbf{Mayor horizonte de entrenamiento.} El número
        máximo de épocas sube de $40$ a $50$ y la paciencia
        de \textit{early stopping} de $7$ a $8$, pues el
        \textit{backbone} más profundo requiere más
        iteraciones para converger.
\end{enumerate}

Bajo este régimen, ambas variantes se entrenan con los mismos
hiperparámetros para que la única variable que distingue a
DeepCNN-BN de DeepCNN-BN-GAP sea, efectivamente, el cabezal
(FC vs GAP).

\subsection{Protocolo de Entrenamiento}
\label{sec:proto}
La implementación se realiza íntegramente en \textbf{Python} con
\textbf{PyTorch}, mediante \textit{scripts} locales organizados
en el directorio \texttt{src/}. El flujo de trabajo reproducible
es el siguiente:

\begin{enumerate}
  \item Crear un entorno virtual aislado:
        \texttt{python -m venv .venv} y activarlo.
  \item Instalar dependencias:
        \texttt{pip install -r requirements.txt}.
  \item Lanzar el entrenamiento de un modelo específico:\\
        \texttt{python src/models/train.py -{}-model N\\
        -{}-data\_dir datasets/Training},
        con \texttt{N} $\in \{1, 2, 3\}$.
  \item Los \textit{checkpoints} de pesos se almacenan en
        \texttt{outputs/results/} y las métricas por época se
        registran incrementalmente en
        \texttt{outputs/results/metrics.csv}.
  \item La evaluación final sobre el conjunto de prueba se
        ejecuta con:\\
        \texttt{python src/models/evaluate.py -{}-model N\\
        -{}-data\_dir datasets/Testing}.
\end{enumerate}

Se aplica una división interna \textbf{80/20} sobre el conjunto
\textit{Training} para obtener un \textit{validation set} usado
en la monitorización de pérdida y métricas durante el
entrenamiento y en la selección del mejor \textit{checkpoint}.
El conjunto \textit{Testing} original se reserva como
\textit{holdout} independiente y solo se emplea en la evaluación
final reportada en la sección de resultados.

\subsection{Hiperparámetros Modelo 3}
\label{sec:hiperparam-m3}

La Tabla~\ref{tab:hiperparam-m3} resume los hiperparámetros
empleados para entrenar el Modelo~3 (MRIResNet).
Las diferencias más relevantes respecto al Modelo~1 son el uso
de \textbf{AdamW} con \textit{weight decay} desacoplado ---en
lugar de Adam sin regularización--- y un \textit{learning rate}
inicial diez veces menor, coherente con la mayor profundidad de la
red y la presencia de \textit{Batch Normalization} en todos los
bloques residuales.

\begin{table}[!t]
\centering
\caption{Hiperparámetros del Modelo~3 (MRIResNet).}
\label{tab:hiperparam-m3}
\renewcommand{\arraystretch}{1.15}
\footnotesize
\begin{tabular}{@{}ll@{}}
\toprule
\textbf{Hiperparámetro} & \textbf{MRIResNet (M3)} \\
\midrule
Optimizador                 & AdamW                     \\
\quad $(\beta_1, \beta_2)$  & $(0{,}9,\;0{,}999)$       \\
Learning rate inicial       & $1\times10^{-4}$          \\
\textit{Weight decay}       & $1\times10^{-4}$          \\
Scheduler                   & ReduceLROnPlateau         \\
\quad factor / patience     & $0{,}5$ / $5$             \\
\quad $\eta_{\min}$         & $1\times10^{-6}$          \\
Función de pérdida          & CrossEntropyLoss          \\
\quad con                   & \texttt{class\_weight}    \\
\textit{Batch size}         & $32$                      \\
\textit{Epochs} máx.        & $50$                      \\
\textit{Early stopping}     & patience $= 10$           \\
\quad criterio              & F1-macro val.             \\
\textit{Split} train/val  & 80\% / 20\% \\
\textit{Seed}               & $42$                      \\
\midrule
Canales de entrada          & $5$ (R, G, B, CLAHE, Sobel) \\
Parámetros entrenables      & $1\,676\,284$             \\
\bottomrule
\end{tabular}
\end{table}

La elección de \textbf{AdamW} frente a Adam responde a que el
\textit{weight decay} desacoplado de AdamW actúa como
regularización $L_2$ explícita independiente del \textit{learning
rate}, lo cual es especialmente beneficioso en redes con
\textit{Batch Normalization}~\cite{loshchilov2019adamw}.
El uso de \texttt{class\_weight} calculado con la fórmula
\textit{balanced} de \texttt{scikit-learn} se mantiene respecto
al Modelo~1 para corregir el desbalance documentado en la
Sección~\ref{sec:hiperparam}.

El protocolo de entrenamiento es análogo al descrito en la
Sección~\ref{sec:proto}: división 80/20, monitorización del
F1-macro de validación para el \textit{scheduler} y el
\textit{early stopping}, y registro de métricas por época en
un archivo \texttt{.csv}.
